# Cervical Cancer Risk Factors - Week 3 Dataset Selection and Final Project Planning Lab

## Purpose and Overview

This notebook is the Dataset Selection and Final Project Planning Lab for the final
course project. It examines the UCI Cervical Cancer (Risk Factors) dataset and evaluates
whether it is suitable for binary classification: predicting whether a patient's cervical
biopsy will be positive based on her age, sexual and reproductive history, contraceptive
use, smoking, and STD history. The practical use is risk stratification, in order to help a clinic
decide which patients to prioritize for follow-up screening. It is a decision-support
idea and does not replace diagnostic testing.

The notebook has four sections:
1. **Data Selection:** the dataset, its source, why I selected it, and a check against the
   project requirements.
2. **Dataset Exploration:** importing the data, structure, target variable, predictors,
   descriptive statistics, and two exploratory visualizations.
3. **Data Quality Assessment:** missing values, duplicates, and impossible values, plus
   anticipated preprocessing steps and potential challenges.
4. **Project Planning:** the prediction problem, target variable, and candidate
   evaluation metrics.

**How to use this notebook:** run every cell from top to bottom
(**Runtime > Restart session and run all** in Colab).

## Dataset Source and Provenance

**Dataset:** Cervical Cancer (Risk Factors)

**Source:** UCI Machine Learning Repository
<https://archive.ics.uci.edu/dataset/383/cervical+cancer+risk+factors>

**Collected at:** Hospital Universitario de Caracas, Caracas, Venezuela (2017).

**License:** CC BY 4.0

**Size:** 858 patients and 36 columns.

**Target:** `Biopsy` (1 = positive, 0 = negative). About 6.4 percent of patients (55 of 858) are
biopsy positive. The other three outcome columns (`Hinselmann`, `Schiller`, `Citology`)
are screening-test results that are strongly related to the biopsy result, so they are
excluded from the predictors.

**Predictors:** 32 columns covering age, sexual and reproductive history, smoking,
contraceptive use, STD history, and prior diagnoses.

**Requirement check:**

| Requirement | Needed | This dataset |
|---|---|---|
| Binary target | Yes | `Biopsy` (0/1) |
| Observations | at least 300 | 858 |
| Predictors | at least 10 | 32 |

**Known data quality issues** (details and planned handling are in the notebook):
- Missing answers are coded as `?` in the raw file (3,622 cells in total; 799 of 858
  patients are missing at least one value).
- Two columns, `STDs: Time since first diagnosis` and `STDs: Time since last diagnosis`,
  are about 92 percent missing.
- There are 23 duplicate rows.
- Two patients have a first-intercourse age that is later than their current age.
- Very few positive cases, so the classes are strongly imbalanced.

**Redistribution:** the CSV (`risk_factors_cervical_cancer.csv`) is included in this
repository under the dataset's CC BY 4.0 license.

## Required Software and Libraries

| Library | Purpose |
|---|---|
| pandas | Loading and summarizing the data |
| numpy | Numerical operations and setting the random seed |
| matplotlib | Plotting |

All libraries are pre-installed in Google Colab. 

A fixed random seed (`RANDOM_SEED = 42`) is set in the Setup cell. 

## Setup / Installation Instructions

**Option A - Google Colab (recommended, no setup required):**
Open the notebook using the Colab link below.

**Colab link:** https://colab.research.google.com/drive/1xLBL_U-zmi-CiAToyAxRE1tTvtoTVmJk#scrollTo=GAbckks39RAq

**Option B - Local Jupyter / VS Code:**
1. Clone this repository:
   ```
   git clone https://github.com/katerina-smith/Programming_for_HDS_Week3.git
   cd Programming_for_HDS_Week3
   ```
   ```
2. Open `Week3_FinalProjectPlanningLab.ipynb`.

## How to Execute the Notebook

The data-loading cell uses a two-step approach so it works in any environment without
editing the code:
1. If `risk_factors_cervical_cancer.csv` is in the current working directory, it is loaded
   from there.
2. Otherwise, it is loaded directly from this project's GitHub repository:
   `https://raw.githubusercontent.com/katerina-smith/Programming_for_HDS_Week3/main/risk_factors_cervical_cancer.csv`

In the raw file, a `?` marks a missing answer. The load step converts these to proper
missing values (`NaN`).

Run every cell from top to bottom:
- In Colab: **Runtime > Restart session and run all**.
- In Jupyter: **Kernel > Restart Kernel and Run All Cells**.

No manual inputs, file uploads, or configuration changes are required.

## Expected Outputs

- A requirements check showing the binary target, number of observations, and number of
  predictors.
- Dataset structure (`head` and `info`), the target distribution, and descriptive
  statistics for the predictors.
- Two figures: the class balance of the target, and the share of patients with a positive
  biopsy with vs. without each risk factor.
- Data quality tables: missing values, duplicate rows, and impossible values.
- Written findings, anticipated preprocessing steps, potential challenges, the proposed
  prediction problem, and candidate evaluation metrics.

## Dataset Requirements

- `risk_factors_cervical_cancer.csv` must be in the same folder as the notebook (included
  in this repository), or it will be retrieved automatically from this project's GitHub
  repository.
- No other files are required.

## Assumptions and Limitations

- **Few positive cases.** Only 55 of 858 patients have a positive biopsy, so results will
  be uncertain.
- **Non-random missing data.** Patients could skip questions for privacy reasons.
- **Single hospital, self-reported data.** Results may not generalize to other
  populations, and sensitive history may be under-reported.
- **Observational data.** Relationships are associations and do not establish causes.

##AI Use

Anthropic. Claude. Used to improve documentation, give examples of optimal 
reproducible workflow, and troubleshoot/optimize code. Used to explore candidate evaluation metrics. 
Also used to generate examples of README documentation in GitHub.
